% =========================================================
% Proof: Differentiability Implies Continuity
% Source: volume-iv/calculus/differentiation/notes/notes-definitions.tex
% Mode: standard
% =========================================================

\subsection*{Differentiability Implies Continuity}
\label{prf:diff-implies-cont}

\begin{remark*}[Return]
\hyperref[thm:diff-implies-cont]{$\leftarrow$ Back to Theorem (Differentiability Implies Continuity) in Notes}
\end{remark*}

\begin{proof}
Let $f$ be defined on an open interval $I$ containing $c$, and assume $f$ is
differentiable at $c$. The claim is that $f$ is continuous at $c$, i.e.,
$\lim_{x \to c} f(x) = f(c)$, equivalently $\lim_{x \to c} [f(x) - f(c)] = 0$.

\ProofStep{1}{Secant identity.}
For any $x \in I$ with $x \neq c$, the vertical displacement admits the factorization
\[
    f(x) - f(c)
    \;=\;
    \underbrace{\left(\frac{f(x) - f(c)}{x - c}\right)}_{\text{difference quotient}}
    \cdot
    \underbrace{(x - c)}_{\Delta x}.
\]
This identity holds on all of $I \setminus \{c\}$ by arithmetic; no limit has
been taken.

\ProofStep{2}{Distribute the limit over the product.}
Taking the limit as $x \to c$ and applying \ByThm{Product Rule for Limits}:
\[
    \lim_{x \to c} [f(x) - f(c)]
    \;=\;
    \left(\lim_{x \to c} \frac{f(x) - f(c)}{x - c}\right)
    \cdot
    \left(\lim_{x \to c} (x - c)\right).
\]
The Product Rule applies because both limits on the right exist — the first by
the differentiability hypothesis, the second trivially.

\ProofStep{3}{Evaluate each factor.}
By \ByDef{Differentiability at a point}, the first factor equals $f'(c)$, a
finite real number. The second factor satisfies $\lim_{x \to c}(x - c) = 0$
by the \ByThm{Identity Limit}. Therefore:
\[
    \lim_{x \to c} [f(x) - f(c)] = f'(c) \cdot 0 = 0.
\]

\ProofStep{4}{Conclude continuity.}
Since $\lim_{x \to c}[f(x) - f(c)] = 0$, it follows that
$\lim_{x \to c} f(x) = f(c)$. By \ByDef{Continuity at a point}, $f$ is
continuous at $c$.
\end{proof}

\begin{remark*}[Proof shape]
The proof hinges on the algebraic factorization of $f(x) - f(c)$ into the
product of the difference quotient and the displacement. This converts the
continuity problem into a limit-of-product problem, where the derivative
hypothesis controls one factor and the displacement annihilates the product.

\FlashProof{1. Factor $f(x)-f(c)$ as (difference quotient)$\cdot(x-c)$. 2. Apply the Product Rule for Limits — both factors have limits. 3. Differentiability gives $f'(c)$ finite; displacement gives $0$. 4. Product is $0$, so $\lim f(x)=f(c)$.}
\end{remark*}

\begin{remark*}[Dependencies]
\begin{itemize}
  \item \hyperref[def:derivative]{Definition of Differentiability at a Point}: Identifies the limit of the difference quotient as $f'(c)$.
  \item \hyperref[prop:limit-product-rule]{Product Rule for Limits}: Distributes the limit over the factored expression.
  \item \hyperref[prop:identity-limit]{Identity Limit}: Establishes $\lim_{x \to c}(x-c) = 0$.
  \item \hyperref[def:continuity]{Definition of Continuity at a Point}: Frames the conclusion.
\end{itemize}
\FlashUses{def:derivative, prop:limit-product-rule, prop:identity-limit, def:continuity}
\end{remark*}